% Part: incompleteness
% Chapter: representability-in-q
% Section: introduction

\documentclass[../../../include/open-logic-section]{subfiles}

\begin{document}

\olfileid[fr]{inc}{req}{int}
\olsection{Introduction}

Les théorèmes d'incomplétude s'appliquent aux théories dans
lesquelles on peut exprimer et démontrer des faits élémentaires sur
les fonctions calculables. Nous décrirons une théorie très faible
de ce type, appelée « $\Th{Q}$ » (ou parfois « arithmétique de Robinson »,
du nom de Raphael Robinson). Nous préciserons ce que signifie,
pour une fonction, être \emph{représentable} dans $\Th{Q}$, puis
nous démontrerons le résultat suivant :
\begin{quote}
  Une fonction totale est représentable dans $\Th{Q}$ si et seulement si
  elle est calculable.
\end{quote}
Cela nous fournit notamment un autre modèle de la calculabilité.
Nous nous en servirons aussi pour montrer que l'ensemble
$\Setabs{!A}{\Th{Q} \Proves !A}$ n'est pas décidable, par réduction
du problème de l'arrêt. Au terme de notre étude, nous aurons
démontré des résultats beaucoup plus forts.

Le langage de $\Th{Q}$ est celui de l'arithmétique; la théorie
$\Th{Q}$ est axiomatisée par les axiomes suivants, auxquels s'ajoutent
les axiomes et règles de la logique du premier ordre avec
!!{identity} :
\begin{align*}
& \lforall[x][\lforall[y][(\eq[x'][y'] \lif \eq[x][y])]] \tag{$!Q_1$}\\
& \lforall[x][\eq/[\Obj 0][x']] \tag{$!Q_2$}\\
& \lforall[x][(\eq[x][\Obj 0] \lor \lexists[y][\eq[x][y']])] \tag{$!Q_3$}\\
& \lforall[x][\eq[(x + \Obj 0)][x]] \tag{$!Q_4$}\\
& \lforall[x][\lforall[y][\eq[(x + y')][(x + y)']]] \tag{$!Q_5$}\\
& \lforall[x][\eq[(x \times \Obj 0)][\Obj 0]] \tag{$!Q_6$}\\
& \lforall[x][\lforall[y][\eq[(x \times y')][((x \times y) + x)]]] \tag{$!Q_7$}\\
& \lforall[x][\lforall[y][(x < y \liff \lexists[z][\eq[(z' + x)][y]])]] \tag{$!Q_8$}
\end{align*}
Pour chaque entier naturel~$n$, appelons \emph{numéral}
$\num{n}$ le terme $\Obj{0}^{\prime\prime\ldots\prime}$ comportant
au total $n$ marques de prime. Ainsi, $\num{0}$ est simplement
la !!{constant}~$\Obj{0}$, $\num{1}$ est $\Obj{0}'$,
$\num{2}$ est $\Obj{0}''$, et ainsi de suite.

En tant que théorie de l'arithmétique, $\Th{Q}$ est
\emph{extrêmement} faible : on ne peut même pas y démontrer des
énoncés aussi simples que $\lforall[x][\eq/[x][x']]$ ou
$\lforall[x][\lforall[y][(x + y) = (y + x)]]$. Nous verrons
pourtant que son intérêt tient en grande partie \emph{à cette
faiblesse même}. En effet, elle possède tout juste assez de force
pour que le théorème d'incomplétude s'y applique. $\Th{Q}$ est
également intéressante parce qu'elle n'a qu'un nombre
\emph{fini} d'axiomes.

On obtient une théorie plus forte que $\Th{Q}$,
l'\emph{arithmétique de Peano} $\Th{PA}$, en ajoutant à~$\Th{Q}$
le schéma d'induction
\[
(!A(\Obj 0) \land \lforall[x][(!A(x) \lif !A(x'))]) \lif \lforall[x][!A(x)]
\]
où $!A(x)$ est une formule quelconque. Si $!A(x)$ contient des
!!{variable}s libres autres que $x$, nous ajoutons en tête des
quantificateurs universels pour les lier toutes, de sorte que
l'instance correspondante du schéma d'induction soit !!a{sentence}.
Par exemple, si $!A(x,y)$ contient aussi la !!{variable} libre~$y$,
l'instance correspondante est
\[
\lforall[y][((!A(\Obj 0) \land \lforall[x][(!A(x) \lif !A(x'))]) \lif
  \lforall[x][!A(x)])]
\]
Les instances du schéma d'induction permettent de démontrer,
à partir des axiomes de~$\Th{PA}$, bien davantage qu'à partir
de ceux de~$\Th{Q}$. En réalité, il faut beaucoup de travail
pour trouver des énoncés « naturels » sur les entiers naturels
qui ne soient pas démontrables dans l'arithmétique de Peano!{}

\begin{defn}
\ollabel{defn:representable-fn}
  Une fonction totale $f\colon\mathbb{N}^{k+1}\to\mathbb{N}$ est dite
  {\em représentable dans $\Th{Q}$}
  s'il existe une formule $!A_f(x_0,\dots,x_k,y)$ telle que, chaque
  fois que $f(n_0,\dots,n_k) = m$, $\Th{Q}$ démontre :
\begin{enumerate}
\item\ollabel{defn:rep:a} $!A_f(\num{n_0}, \dots, \num{n_k}, \num{m})$
\item\ollabel{defn:rep:b} $\lforall[y][(!A_f(\num{n_0}, \dots,
\num{n_k}, y) \lif \num{m} = y)]$.
\end{enumerate}
\end{defn}

On peut formuler cette définition autrement. Par exemple, il
reviendrait au même d'exiger que $\Th{Q}$ démontre
$\lforall[y][(!A_f(\num{n_0}, \dots,
\num{n_k}, y) \liff \eq[y][\num{m}])]$.\footnote{La représentabilité exige ces preuves pour chaque tuple de numéraux fixé. Elle ne demande pas une preuve dans $\Th{Q}$ de la totalité et de l'unicité avec des entrées universellement quantifiées.}

\begin{thm}
\ollabel{thm:representable-iff-comp}
Une fonction totale est représentable dans $\Th{Q}$ si et seulement si
elle est calculable.
\end{thm}

La démonstration comporte deux implications. Celle de gauche à
droite est assez directe une fois la syntaxe arithmétisée. L'autre
demande davantage de travail. En voici l'idée : nous choisissons
la notion de fonction « récursive générale » pour préciser celle
de fonction « calculable », puis nous montrons que toute fonction
récursive générale est représentable dans~$\Th{Q}$. Rappelons
qu'une fonction est récursive générale si elle peut être définie
à partir de la fonction nulle~$\Zero$, de la fonction successeur
$\Succ$ et des projections~$\Proj{n}{i}$, au moyen de la
composition, de la récurrence primitive et de la minimisation
régulière, c'est-à-dire de la recherche non bornée appliquée à
une fonction régulière. Pour montrer que toute fonction récursive
générale est représentable dans~$\Th{Q}$, nous pourrions donc
établir que les fonctions de base sont représentables et que la
composition, la récurrence primitive et la minimisation régulière
préservent la représentabilité. Autrement dit, il suffirait de
montrer que la classe des fonctions représentables est
\emph{close par} ces trois opérations. Toute fonction récursive
générale serait alors représentable.

Or, il est difficile de montrer que les fonctions représentables
sont closes par récurrence primitive. Pour éviter cette étape,
nous montrons d'abord qu'on peut en fait se passer de la
récurrence primitive : toute fonction récursive générale peut
être définie à partir de fonctions de base par composition et
minimisation régulière seulement. Pour cela, nous montrerons
qu'une minimisation régulière particulière permet de réaliser
la récurrence primitive. Il faut toutefois ajouter aux fonctions
de base l'addition, la multiplication et la fonction
caractéristique de la relation d'égalité~$\Char{=}$. Nous
pourrons alors démontrer le théorème en établissant que toutes
\emph{ces} fonctions de base sont représentables dans~$\Th{Q}$
et que la classe des fonctions représentables est close par
composition et par minimisation régulière.

\end{document}
